% !TeX root = ../tesis.tex

% Thesis Abstract -----------------------------------------------------

%\begin{abstractslong}    %uncommenting this line, gives a different abstract heading
\begin{abstracts}        %this creates the heading for the abstract page
\addcontentsline{toc}{chapter}{\protect\numberline{}Resumen}

Los eritrocitos presentan propiedades ópticas que dependen de su geometría, tamaño, composición y orientación, por lo que alteraciones morfológicas y variaciones en la concentración de hemoglobina pueden modificar la manera en que absorben y esparcen la luz. Esta sensibilidad ha motivado el uso de técnicas ópticas, como la citometría de flujo, para la caracterización y diferenciación celular. Sin embargo, gran parte de los estudios numéricos del esparcimiento de luz por eritrocitos se han concentrado en longitudes de onda cercanas a 632.8 nm, fuera de las principales bandas de absorción de la hemoglobina. En este trabajo de tesis de licenciatura se estudian numéricamente las propiedades ópticas de eritrocitos sanos y de dos alteraciones morfológicas representativas, el esferocito y el microcito, en el intervalo espectral de 400 a 600 nm. Las distintas morfologías se modelaron mediante óvalos de Cassini, mientras que sus propiedades ópticas se incorporaron a través de funciones dieléctricas dependientes de la longitud de onda y de la concentración de hemoglobina, obtenidas a partir de datos experimentales mediante las relaciones de Kramers-Kronig sustractivas. Las simulaciones se realizaron mediante el método \textit{Finite Integration Technique} en CST Studio Suite. Los parámetros óptimos para realizar las simulaciones se establecieron mediante un análisis paramétrico y comparación con la solución analítica de Mie para una esfera, obteniéndose errores relativos inferiores al 5\%, y posteriormente se contrastó el patrón angular de un eritrocito con resultados numéricos reportados en la literatura. Los resultados muestran que la respuesta óptica depende del efecto combinado de la geometría y de las propiedades dieléctricas. La banda de Soret, alrededor de 431 nm, presenta las mayores diferencias en la eficiencia de extinción entre las morfologías consideradas. Aunque el esparcimiento se concentra predominantemente en la dirección frontal, las mayores diferencias relativas aparecen en regiones laterales y posteriores. Finalmente, el análisis mediante ventanas angulares permitió identificar una región alrededor de $\theta\approx130^\circ$ en la que el contraste entre eritrocitos sanos y patológicos aumenta de manera consistente para las tres longitudes de onda analizadas. Estos resultados sugieren condiciones espectrales y angulares potencialmente útiles para futuros esquemas de diferenciación celular mediante citometría de flujo.



\end{abstracts}
%\end{abstractlongs}


% ----------------------------------------------------------------------